% 第 6 章
\bichapter{芯片收尾与可制造性设计}{Chip Finishing and Design for Manufacturing}
\label{chap:dfm}

Fusion Compiler 工具提供芯片收尾（chip finishing）以及可制造性设计（DFM）与良率（yield）相关能力，可在设计流程各阶段应用，以应对芯片制造中的工艺设计问题。

有关芯片收尾与可制造性设计功能，请参阅以下主题：
\begin{itemize}
  \item 插入 Tap Cell（Inserting Tap Cells）
  \item 执行 Boundary Cell 插入（Performing Boundary Cell Insertion）
  \item 查找并修复天线违例（Finding and Fixing Antenna Violations）
  \item 插入冗余 Via（Inserting Redundant Vias）
  \item 优化线长与 Via 数量（Optimizing Wire Length and Via Count）
  \item 减小临界面积（Reducing Critical Areas）
  \item 插入金属--绝缘体--金属电容（Inserting Metal-Insulator-Metal Capacitors）
  \item 插入 Filler Cell（Inserting Filler Cells）
  \item 插入 Metal Fill（Inserting Metal Fill）
\end{itemize}

% ============================================================
\section{插入 Tap Cell}
\subsection*{Inserting Tap Cells}

Tap cell 是带有 well tie、substrate tie 或二者兼有的特殊非逻辑单元。当库中多数或全部标准单元不含 substrate/well tap 时，通常需要使用此类单元。一般设计规则会规定标准单元内每个晶体管到 well 或 substrate tap 的最大允许距离。

在全局布局之前（floorplanning 阶段），可在 block 中插入 tap cell，形成二维阵列结构，以确保后续放置的标准单元均满足最大 diffusion-to-tap 距离限制。插入后应目视检查，确认所有标准单元可布局区域均已被 tap cell 妥善保护。

Fusion Compiler 提供两种插入 tap cell 阵列的方法：
\begin{itemize}
  \item 基于 site row 以及 tap cell 之间的最大距离插入 tap cell 阵列\\
        使用 \cmd{create\_tap\_cells} 命令，见「使用 \cmd{create\_tap\_cells} 命令」。
  \item 使用指定的 offset 与 pitch 插入 tap cell 阵列，并可选择在边界附近额外插入 tap cell\\
        使用 \cmd{create\_cell\_array} 命令，见「使用 \cmd{create\_cell\_array} 命令」。
\end{itemize}

先进工艺节点常需额外插入 tap cell，以管理 substrate 与 well 噪声。在标准 tap cell 插入之后，还可使用以下能力：
\begin{itemize}
  \item 插入 tap wall\\
        Tap wall 是一行或一列无间隙线性排列的 tap cell。可在 block、macro 或 hard placement blockage 边界之外或之内插入。边界外的称为 exterior tap wall；边界内的称为 interior tap wall。
        \begin{itemize}
          \item 插入 exterior tap wall：使用 \cmd{create\_exterior\_tap\_walls}，见「插入 Exterior Tap Wall」。
          \item 插入 interior tap wall：使用 \cmd{create\_interior\_tap\_walls}，见「插入 Interior Tap Wall」。
        \end{itemize}
  \item 插入 tap mesh\\
        Tap mesh 是二维 tap cell 阵列，每个 mesh window 中有一个 tap cell。使用 \cmd{create\_tap\_meshes}，见「插入 Tap Mesh」。
  \item 插入 dense tap array\\
        Dense tap array 的 tap 间距小于 \cmd{create\_tap\_cells} 所用间距。使用 \cmd{create\_dense\_tap\_cells}，见「插入 Dense Tap Array」。
\end{itemize}

\subsection{使用 \cmd{create\_tap\_cells} 命令}
\subsubsection*{Using the create\_tap\_cells Command}

要基于 site row 添加 tap cell 阵列，使用 \cmd{create\_tap\_cells}。必须指定用于插入的库单元名（\opt{-lib\_cell}）以及 tap cell 之间的最大距离（微米，\opt{-distance}）。

例如：
\begin{lstlisting}
fc_shell> create_tap_cells -lib_cell myreflib/mytapcell -distance 30
\end{lstlisting}

\begin{noteBox}
若应用选项 \appopt{place.legalize.enable\_pin\_color\_alignment\_check} 为 \cmd{true}（默认 \cmd{false}），命令在插入时会确保 tap cell 内部 pin 与金属对齐到相应颜色的 routing track，并移动 tap cell 以避免颜色对齐违例。为避免因移动导致 signoff DRC 错误，在该应用选项为 \cmd{true} 时应略微减小 tap 距离。
\end{noteBox}

默认情况下，\cmd{create\_tap\_cells} 在整个 block 的每一行插入 tap cell。工具从行左边缘开始，按指定 tap 距离确定后续位置。此外，若在相邻于 block 边界、hard macro 或 hard placement blockage 的每个行边缘的最小 tap 距离（指定距离的一半）内不存在 tap cell，工具会额外插入一个 tap cell。图~\ref{fig:default-tap-placement} 给出采用 every-row 插入模式时的默认布局；注意 hard macro 右侧会额外加入 tap cell，以确保距 hard macro 边缘在最小 tap 距离内存在 tap cell。

\begin{figure}[htbp]
\centering
\fbox{\parbox{0.85\textwidth}{\centering\vspace{1.5cm}
\textit{（原书 Figure 85：Default Tap Cell Placement）}\\[0.3em]
\small 请参阅原版 PDF 插图；可将截图放入 \texttt{figures/} 目录后用 \texttt{\textbackslash includegraphics} 替换。
\vspace{1.5cm}}}
\caption{默认 Tap Cell 布局}
\label{fig:default-tap-placement}
\end{figure}

可修改以下默认行为方面：
\begin{itemize}
  \item \textbf{镜像（MX）行所用库单元}\\
        使用 \opt{-mirrored\_row\_lib\_cell} 指定镜像行所用库单元。
  \item \textbf{插入图案}\\
        使用 \opt{-pattern} 指定以下之一：
        \begin{itemize}
          \item \texttt{every\_row}（默认）：每行都插入。此时 \opt{-distance} 应约为工艺设计规则中最大 diffusion-to-tap 值的两倍。
          \item \texttt{every\_other\_row}：仅在奇数行插入。\opt{-distance} 应约为最大 diffusion-to-tap 值的两倍。
          \item \texttt{stagger}：每行都插入，偶数行相对奇数行偏移半个 \opt{-offset}，形成类似棋盘的图案。\opt{-distance} 应约为最大 diffusion-to-tap 值的四倍。
        \end{itemize}
  \item \textbf{相对行左边缘的偏移}\\
        使用 \opt{-offset} 将图案起点向右平移指定微米距离。
  \item \textbf{固定单元的处理}\\
        默认在计算出的 tap 位置插入，不论该处是否被固定单元占用。使用 \opt{-skip\_fixed\_cells} 可避免在固定单元位置插入。默认将该固定单元视为 blockage 并在该处打断行，从而可能在固定单元两侧各插入一个 tap cell。若要不打断行又避免重叠，可同时使用 \opt{-preserve\_distance\_continuity} 与 \opt{-skip\_fixed\_cells}：命令会平移 tap cell 以避免重叠；\opt{-preserve\_distance\_continuity} 仅平移那些计算位置被所列库单元实例占用的 tap cell。
\end{itemize}

例如，下列命令会平移计算位置被 \texttt{mylib/cell1} 固定实例占用的 tap cell：
\begin{lstlisting}
fc_shell> create_tap_cells -lib_cell mylib/tap -distance 30 \
   -skip_fixed_cells -preserve_distance_continuity mylib/cell1
\end{lstlisting}

使用 \opt{-skip\_fixed\_cells} 时，还可用 \opt{-min\_horizontal\_periphery\_spacing} 指定边界与位于 blockage 上方或下方的角部 tap cell 之间的最小水平间距（以 unit tile 数为单位）。图~\ref{fig:tap-boundary-spacing} 示出有无 boundary cell 时该选项对布局的影响。

\begin{figure}[htbp]
\centering
\fbox{\parbox{0.85\textwidth}{\centering\vspace{1.5cm}
\textit{（原书 Figure 86：Tap Cell Boundary Row Spacing）}\\[0.3em]
\small 请参阅原版 PDF 插图。
\vspace{1.5cm}}}
\caption{Tap Cell 边界行间距}
\label{fig:tap-boundary-spacing}
\end{figure}

\begin{itemize}
  \item \textbf{额外 tap cell 的添加}\\
        可用下列一种或两种方法阻止额外插入：
        \begin{itemize}
          \item \opt{-row\_end\_tap\_bypass}：声明 boundary cell 已含 tap，故边界不再额外插入。注意命令不会验证 block 是否确有 boundary cell，或其是否真含 tap。
          \item \opt{-at\_distance\_only}：仅允许在指定 tap 距离、半距离或四分之一距离（仅 stagger）处插入。注意使用该选项可能导致 DRC 违例。
        \end{itemize}
  \item \textbf{插入区域}\\
        使用 \opt{-voltage\_area} 将插入限制在指定 voltage area。
  \item \textbf{命名约定}\\
        默认命名：\texttt{tapfiller!library\_cell\_name!number}。\\
        用 \opt{-separator} 可更改默认分隔符 ``!''。\\
        用 \opt{-prefix} 可为特定运行加前缀，命名变为：\texttt{tapfiller!prefix!library\_cell\_name!number}。
\end{itemize}

\subsection{使用 \cmd{create\_cell\_array} 命令}
\subsubsection*{Using the create\_cell\_array Command}

要按特定布局添加 tap cell 阵列，使用 \cmd{create\_cell\_array}。必须指定库单元名（\opt{-lib\_cell}）以及 $x$、$y$ 方向的 pitch。

例如：
\begin{lstlisting}
fc_shell> create_cell_array -lib_cell myreflib/mytapcell \
   -x_pitch 10 -y_pitch 10
\end{lstlisting}

默认在整个 block 的指定网格每个位置插入 tap cell。工具从 core 区域左下角开始，按 $x$/$y$ pitch 确定后续位置。插入的 tap cell 会 snap 到 site row，并以该行默认方向放置。若位置被固定单元、hard macro、soft macro 或 power-switch cell 占用，则不在该处插入。

可修改以下默认行为：
\begin{itemize}
  \item \textbf{插入区域}\\
        \opt{-voltage\_area}：限制到指定 voltage area。\\
        \opt{-boundary}：限制到指定区域。\\
        若两者同时使用，仅在重叠区域插入。
  \item \textbf{相对插入区域边界的偏移}\\
        \opt{-x\_offset}：起点向右平移（微米）。\\
        \opt{-y\_offset}：起点向上平移（微米）。
  \item \textbf{插入图案}\\
        \opt{-checkerboard}：在网格每隔一个位置插入。
        \begin{itemize}
          \item \opt{-checkerboard even}：左下角放置 tap cell。
          \item \opt{-checkerboard odd}：左下角留空。
        \end{itemize}
        某些设计方法要求在棋盘图案中位于特定单元（如 critical area breaker 或 boundary cell）上方或下方的空位也插入 tap cell。可用 \opt{-preserve\_boundary\_row\_lib\_cells} 指定相关库单元以启用这些额外插入。
  \item \textbf{Snap 行为}\\
        设置 \opt{-snap\_to\_site\_row false} 可禁用 snap 到 site row；此时默认方向为 R0。
  \item \textbf{方向}\\
        用 \opt{-orient} 指定：R0、R90、R180、R270、MX、MY、MXR90 或 MYR90。
  \item \textbf{命名约定}\\
        默认：\texttt{headerfooter\_\_library\_cell\_name\_R\#\_C\#\_number}。\\
        使用 \opt{-prefix} 时：\texttt{prefix\_\_library\_cell\_name\_R\#\_C\#\_number}。
\end{itemize}

\subsection{插入 Exterior Tap Wall}
\subsubsection*{Inserting Exterior Tap Walls}

Exterior tap wall 是放置在 block、macro 或 hard placement blockage 边界之外的一行或一列无间隙 tap cell。

使用 \cmd{create\_exterior\_tap\_walls}，必须指定：
\begin{itemize}
  \item \textbf{Tap wall 所用 tap cell}（\opt{-lib\_cell}）\\
        默认以 R0 方向插入；可用 \opt{-orientation} 指定其他方向。\\
        若设计有 FinFET grid，选择 tap cell 时请遵循：
        \begin{itemize}
          \item 水平 tap wall：指定方向下的单元宽度须为 FinFET grid $x$-pitch 的整数倍；
          \item 垂直 tap wall：指定方向下的单元高度须为 FinFET grid $y$-pitch 的整数倍。
        \end{itemize}
        默认命名：\texttt{extTapWall\_\_library\_cell\_name\_R\#\_C\#\_number}。\\
        使用 \opt{-prefix} 时：\texttt{prefix\_\_library\_cell\_name\_R\#\_C\#\_number}。
  \item \textbf{插入侧}（\opt{-side}）\\
        每次只能指定一侧：\texttt{top}、\texttt{bottom}、\texttt{right} 或 \texttt{left}。多侧需多次运行。
  \item \textbf{插入区域}（\opt{-bbox}）\\
        指定包围盒左下与右上角：
\begin{lstlisting}
-bbox {{llx lly} {urx ury}}
\end{lstlisting}
        包围盒必须覆盖指定侧上对象边缘的一部分。若覆盖该侧多条边，则沿最长边插入；若未覆盖任何边，则报错。
\end{itemize}

默认情况下，水平 tap wall 从包围盒左边缘开始并贴齐顶/底边界；垂直 tap wall 从包围盒底边缘开始并贴齐左/右边界。可用 \opt{-x\_margin} 与 \opt{-y\_margin} 指定边距。若含边距后的区域过小，则不插入。应确保插入区域不含 placement blockage、固定单元或 macro，否则可能导致 tap wall 不完整或缺失。

插入时命令不遵循标准单元 row/site，也不跨越对象边界。插入后会固定所插 tap cell 的布局。

\subsection{插入 Interior Tap Wall}
\subsubsection*{Inserting Interior Tap Walls}

Interior tap wall 是放置在 block、macro、hard placement blockage 或非默认 voltage area 的标准单元布局区域边界之内的一行或一列 tap cell。

\begin{noteBox}
仅可对具有单一 site 定义的设计插入 interior tap wall。
\end{noteBox}

使用 \cmd{create\_interior\_tap\_walls}，必须指定：
\begin{itemize}
  \item \textbf{Tap cell}（\opt{-lib\_cell}）\\
        须为与行高、site 宽匹配的单高（single-height）单元。\\
        默认按行方向插入；可用 \opt{-orientation} 更改。\\
        默认命名：\texttt{tapfiller\_\_library\_cell\_name\_R\#\_C\#\_number}。\\
        使用 \opt{-prefix} 时：\texttt{prefix\_\_library\_cell\_name\_R\#\_C\#\_number}。
  \item \textbf{插入侧}（\opt{-side}）\\
        每次仅一侧；多侧需多次运行。
\end{itemize}

默认沿标准单元可布局区域指定侧的所有边无间隙放置。
\begin{itemize}
  \item 用 \opt{-bbox} 指定插入区域包围盒。命令沿包围盒所涵盖可布局区域在指定侧的最长边插入。对齐取决于该边段所含角点数量：
        \begin{itemize}
          \item 无角点：水平 wall 从左边、垂直 wall 从底边开始；
          \item 一个角点：从该角点开始；
          \item 指定侧上两个角点：水平从左边、垂直从底边开始。对含两角的水平 wall，末两单元间距可能缩短以对齐右角；空间不足时末单元可能不贴齐右角。
        \end{itemize}
  \item 对水平 tap wall，可用 \opt{-x\_spacing} 指定单元间隙。若小于单元宽度则间距为 0；若非 site 宽整数倍则向下取整到 site 宽倍数。
\end{itemize}

插入时遵循标准单元 row 与 site；未定义则不插入。Placement blockage、固定单元与 macro 等障碍可能导致不完整或缺失的 tap wall。插入后固定布局。

\subsection{插入 Tap Mesh}
\subsubsection*{Inserting Tap Meshes}

Tap mesh 是每个 mesh window 含一个 tap cell 的二维阵列。通常插在 P\&R block 之外，也可插在完全位于 core 区域内或完全位于 core 区域外的任意区域。

使用 \cmd{create\_tap\_meshes}，必须指定：
\begin{itemize}
  \item \textbf{Tap cell}（\opt{-lib\_cell}）\\
        默认 R0；可用 \opt{-orientation}。\\
        默认命名：\texttt{headerfooter\_\_library\_cell\_name\_R\#\_C\#\_number}。\\
        使用 \opt{-prefix} 时：\texttt{prefix\_\_library\_cell\_name\_R\#\_C\#\_number}。
  \item \textbf{插入区域}（\opt{-bbox}）：\texttt{\{\{llx lly\} \{urx ury\}\}}
  \item \textbf{Mesh window}（\opt{-mesh\_window}）\\
        以 $x$、$y$ 方向 pitch 的整数值指定窗口大小：\texttt{-mesh\_window \{x\_pitch y\_pitch\}}。\\
        Mesh window 与插入区域左下角对齐。
\end{itemize}

插入时不遵循标准单元 row/site。在每个 mesh window 左下角放置 tap cell；若被阻挡，则取到该左下角曼哈顿距离最近的可用位置；若无可用位置则警告并跳过该窗口。插入后固定布局。

\subsection{插入 Dense Tap Array}
\subsubsection*{Inserting Dense Tap Arrays}

Dense tap array 的 tap 间距小于 \cmd{create\_tap\_cells} 所用间距。

使用 \cmd{create\_dense\_tap\_cells}，必须指定：
\begin{itemize}
  \item \textbf{Tap cell}（\opt{-lib\_cell}）\\
        须与运行 \cmd{create\_tap\_cells} 时相同。\\
        默认命名：\texttt{tapfiller!library\_cell\_name!number}；可用 \opt{-separator}、\opt{-prefix}。
  \item \textbf{Tap 距离}（\opt{-distance}）\\
        必须小于 \cmd{create\_tap\_cells} 所用距离。
  \item \textbf{插入区域}（\opt{-bbox}）
\end{itemize}

命令首先移除完全位于插入区域内的既有 tap cell，再按指定距离重新插入以形成更密阵列。为避免 tap 规则违例，应使用与 \cmd{create\_tap\_cells} 相同的插入图案。支持 \opt{-offset}、\opt{-pattern}、\opt{-skip\_fixed\_cells}、\opt{-at\_distance\_only}（详见 man 页）。插入后固定布局。

% ============================================================
\section{执行 Boundary Cell 插入}
\subsection*{Performing Boundary Cell Insertion}

在放置标准单元之前，可为 block 添加 boundary cell。Boundary cell 包括：加在单元行两端以及 core 区域、hard macro、blockage、voltage area 等对象边界周围的 end-cap cell；以及填充水平与垂直 end-cap 之间空隙的 corner cell。End-cap cell 通常为非逻辑单元（如电源轨去耦电容）。因工具接受任意标准单元作为 end-cap，须确保指定合适的单元。

插入 boundary cell 的步骤：
\begin{enumerate}
  \item 用 \cmd{set\_boundary\_cell\_rules} 指定插入要求（见「指定 Boundary Cell 插入要求」）。
  \item 用 \cmd{compile\_boundary\_cells} 或 \cmd{compile\_targeted\_boundary\_cells} 按规则插入（见「插入 Boundary Cell」）。
  \item 用 \cmd{check\_boundary\_cells} 或 \cmd{check\_targeted\_boundary\_cells} 验证布局（见「验证 Boundary Cell 布局」）。
\end{enumerate}

每次运行仅支持单一 site 定义。\cmd{compile\_boundary\_cells} 与 \cmd{compile\_targeted\_boundary\_cells} 根据 \cmd{set\_boundary\_cell\_rules} 所指定库单元确定 site，并仅在使用该 site 定义的区域插入。若设计有多种 site 定义，须对每种各运行一次该流程。

\subsection{指定 Boundary Cell 插入要求}
\subsubsection*{Specifying the Boundary Cell Insertion Requirements}

使用 \cmd{set\_boundary\_cell\_rules} 指定：
\begin{itemize}
  \item 用于 boundary cell 的库单元（见「指定 Boundary Cell 插入所用库单元」）；
  \item 放置规则（见「指定 Boundary Cell 放置规则」）；
  \item 插入单元的命名约定（见「指定 Boundary Cell 命名约定」）；
  \item 创建 routing guide 以遵守 metal cut allowed/forbidden preferred grid extension 规则（见「Boundary Cell 插入期间创建 Routing Guide」）；
\begin{noteBox}
该功能仅由 \cmd{compile\_boundary\_cells} 支持，\cmd{compile\_targeted\_boundary\_cells} 不支持。
\end{noteBox}
  \item 创建 placement blockage 以遵守 minimum jog / minimum separation 规则（见「Boundary Cell 插入期间创建 Placement Blockage」）。
\begin{noteBox}
该功能仅由 \cmd{compile\_boundary\_cells} 支持，\cmd{compile\_targeted\_boundary\_cells} 不支持。
\end{noteBox}
\end{itemize}

\begin{noteBox}
\cmd{set\_boundary\_cell\_rules} 指定的设置保存在设计库中。
\end{noteBox}

\begin{seeAlsoBox}
报告 Boundary Cell 插入要求：见「报告 Boundary Cell 插入要求」。
\end{seeAlsoBox}

\subsubsection{指定 Boundary Cell 插入所用库单元}
\subsubsection*{Specifying the Library Cells for Boundary Cell Insertion}

Boundary cell 包括左右上下边界的 end-cap，以及内/外角单元。可用 \cmd{set\_boundary\_cell\_rules} 的下列选项为各类型指定不同库单元：
\begin{itemize}
  \item \opt{-left\_boundary\_cell}、\opt{-right\_boundary\_cell}：分别指定左、右边界 end-cap 所用的单个库单元。
  \item \opt{-top\_boundary\_cells}、\opt{-bottom\_boundary\_cells}：分别指定顶、底边界 end-cap 的库单元列表。按指定顺序插入；若剩余空间小于当前单元，则插入列表中下一个能放入剩余空间的单元。\\
        对垂直行 block，行自底向顶，故顶边界沿 block 左侧、底边界沿右侧。\\
        对 double-back block 中的翻转行，顶边界单元用于底边界，底边界单元用于顶边界。
  \item \opt{-top\_tap\_cell}、\opt{-bottom\_tap\_cell}：分别指定顶、底边界行上的 tap cell，以确保端 cap 满足最大 diffusion-to-tap 限制。
  \item \opt{-top\_left\_outside\_corner\_cell}、\opt{-top\_right\_outside\_corner\_cell}、\opt{-bottom\_left\_outside\_corner\_cell}、\opt{-bottom\_right\_outside\_corner\_cell}：各外角位置的单个库单元。
  \item \opt{-top\_left\_inside\_corner\_cells}、\opt{-top\_right\_inside\_corner\_cells}、\opt{-bottom\_left\_inside\_corner\_cells}、\opt{-bottom\_right\_inside\_corner\_cells}：各内角位置的库单元列表。工具插入第一个与内角尺寸匹配的角单元；若无一精确匹配，则插入第一个可不违规则放置的单元。
\end{itemize}

图~\ref{fig:boundary-cell-locations} 示出含两个 hard macro/blockage 的水平行 block 上 boundary 与 corner 单元位置；行翻转时位置也会翻转。

\begin{figure}[htbp]
\centering
\fbox{\parbox{0.85\textwidth}{\centering\vspace{1.5cm}
\textit{（原书 Figure 87：Boundary Cell Locations）}\\[0.3em]
\small 请参阅原版 PDF 插图。
\vspace{1.5cm}}}
\caption{Boundary Cell 位置}
\label{fig:boundary-cell-locations}
\end{figure}

\subsubsection{指定 Boundary Cell 放置规则}
\subsubsection*{Specifying Boundary Cell Placement Rules}

默认：
\begin{itemize}
  \item 运行 \cmd{compile\_boundary\_cells} 时，工具以默认方向在 core 区域、hard macro 与 hard placement blockage 周围放置 boundary cell；
  \item 运行 \cmd{compile\_targeted\_boundary\_cells} 时，工具以默认方向在指定对象周围放置。
\end{itemize}

要翻转边界单元方向，可在 \cmd{set\_boundary\_cell\_rules} 中使用一个或多个选项：\opt{-mirror\_left\_boundary\_cell}、\opt{-mirror\_right\_boundary\_cell}、\opt{-mirror\_left\_outside\_corner\_cell}、\opt{-mirror\_right\_outside\_corner\_cell}、\opt{-mirror\_left\_inside\_corner\_cell}、\opt{-mirror\_right\_inside\_corner\_cell}、\opt{-mirror\_left\_inside\_horizontal\_abutment\_cell}、\opt{-mirror\_right\_inside\_horizontal\_abutment\_cell}。不能翻转顶/底边界单元方向。

使用 \cmd{compile\_boundary\_cells} 时，还可修改：
\begin{itemize}
  \item \textbf{翻转行上顶/底内角单元互换}\\
        \opt{-do\_not\_swap\_top\_and\_bottom\_inside\_corner\_cell}：阻止在翻转行顶内角使用底内角单元、在底内角使用顶内角单元。
  \item \textbf{考虑 voltage area}\\
        \opt{-at\_va\_boundary}：在 voltage area 边界两侧插入水平 boundary cell。
  \item \textbf{一单位 tile 间隙}\\
        \opt{-no\_1x}：防止插入产生一单位 tile 间隙。当行长等于两倍角单元宽加一单位 tile 宽时不插入；若等于两倍角单元宽则仍插入。
  \item \textbf{Tap cell 间距}\\
        \opt{-tap\_distance}：顶/底边界行上 tap cell 间距（微米）。
  \item \textbf{短行上的插入}\\
        \opt{-min\_row\_width}：行宽小于该阈值时不插入。
  \item \textbf{子 block 中的插入}\\
        \opt{-insert\_into\_blocks}：递归在子 block 中插入（默认不插入）。
\end{itemize}

\subsubsection{指定 Boundary Cell 命名约定}
\subsubsection*{Specifying the Naming Convention for Boundary Cells}

默认命名：\texttt{boundarycell!library\_cell\_name!number}。\\
用 \opt{-separator} 更改分隔符；用 \opt{-prefix} 加前缀后为：\texttt{boundarycell!prefix!library\_cell\_name!number}。

\subsection{Boundary Cell 插入期间创建 Routing Guide}
\subsubsection*{Creating Routing Guides During Boundary Cell Insertion}

若 technology file 定义了 metal cut allowed 与 forbidden preferred grid extension 规则，可用 \opt{-add\_metal\_cut\_allowed} 在插入期间创建相应 routing guide。此时 \cmd{compile\_boundary\_cells} 会创建：
\begin{itemize}
  \item Metal cut allowed routing guide：覆盖全部可布局 site row 所占区域，再按垂直收缩因子缩小（为最小 site row 高度的 50\%）；
  \item Forbidden preferred grid extension routing guide：覆盖直至 block 边界的剩余区域。
\end{itemize}

\begin{noteBox}
仅 \cmd{compile\_boundary\_cells} 支持；\cmd{compile\_targeted\_boundary\_cells} 不支持。
\end{noteBox}

\subsection{Boundary Cell 插入期间创建 Placement Blockage}
\subsubsection*{Creating Placement Blockages During Boundary Cell Insertion}

若工艺有 minimum jog 或 minimum separation 要求，可在 \cmd{set\_boundary\_cell\_rules} 中配合下列选项，由 \cmd{compile\_boundary\_cells} 创建 placement blockage：
\begin{itemize}
  \item \opt{-min\_horizontal\_jog}：若 macro（含 hard keepout）、hard placement blockage 或 voltage area（含 guard band）的水平边小于指定值，则用 \cmd{compile\_boundary\_cells -add\_placement\_blockage} 创建 blockage 以防违例。
  \item \opt{-min\_vertical\_jog}：对垂直边同理。
  \item \opt{-min\_horizontal\_separation}：若连续水平布局区域小于指定值，则创建 blockage。
  \item \opt{-min\_vertical\_separation}：对垂直方向同理。
\end{itemize}

\begin{noteBox}
默认 \cmd{compile\_boundary\_cells} 会检查指定规则但不创建 blockage；须使用 \opt{-add\_placement\_blockage} 才会创建。该功能仅 \cmd{compile\_boundary\_cells} 支持。
\end{noteBox}

\subsection{报告 Boundary Cell 插入要求}
\subsubsection*{Reporting the Boundary Cell Insertion Requirements}

使用 \cmd{report\_boundary\_cell\_rules} 报告由 \cmd{set\_boundary\_cell\_rules} 指定的要求。仅报告用户指定设置，不报告默认值。

\subsection{移除 Boundary Cell 插入要求}
\subsubsection*{Removing Boundary Cell Insertion Requirements}

使用 \cmd{remove\_boundary\_cell\_rules} 移除一项或多项由 \cmd{set\_boundary\_cell\_rules} 指定的要求。

\subsection{插入 Boundary Cell}
\subsubsection*{Inserting Boundary Cells}

按插入目标选择命令：
\begin{itemize}
  \item 在 core 区域、hard macro 与 placement blockage 周围插入：\cmd{compile\_boundary\_cells}；
  \item 仅在一个或多个 voltage area 周围：\cmd{compile\_boundary\_cells -voltage\_area}；
  \item 仅在特定对象周围：\cmd{compile\_targeted\_boundary\_cells -target\_objects}（可针对特定 macro、voltage area、voltage area shape、placement/routing blockage 与 core 区域）。
\end{itemize}

\cmd{set\_boundary\_cell\_rules} 同时配置两种 compile 命令；部分规则仅被 \cmd{compile\_boundary\_cells} 遵守，详见「指定 Boundary Cell 插入要求」。

\subsection{验证 Boundary Cell 布局}
\subsubsection*{Verifying the Boundary Cell Placement}

用 \cmd{compile\_boundary\_cells} 插入后，用 \cmd{check\_boundary\_cells} 验证。用 \opt{-error\_view} 可生成供 error browser 查看的错误数据文件。该命令检查：
\begin{itemize}
  \item \textbf{缺失的 boundary 或 corner cell}\\
        验证整个边界周围有 boundary/corner cell 且无间隙。图~\ref{fig:check-missing-boundary} 示出有效布局及因缺失单元导致的错误。
\begin{noteBox}
对成熟节点，boundary cell 可能仅插在左、右侧；此时命令只验证这两侧无间隙。
\end{noteBox}
  \item \textbf{不正确的 boundary/corner cell}\\
        验证所用库单元与 \cmd{set\_boundary\_cell\_rules} 指定一致。
  \item \textbf{不正确的方向}\\
        验证各 boundary cell 方向符合允许方向。
  \item \textbf{短行与短边}\\
        验证每行 boundary cell 宽于 \opt{-min\_row\_width}，且各 blockage 水平边大于 \opt{-min\_horizontal\_jog}。
\end{itemize}

\begin{figure}[htbp]
\centering
\fbox{\parbox{0.85\textwidth}{\centering\vspace{1.5cm}
\textit{（原书 Figure 88：Check for Missing Boundary or Corner Cells）}\\[0.3em]
\small 请参阅原版 PDF 插图。
\vspace{1.5cm}}}
\caption{检查缺失的 Boundary/Corner Cell}
\label{fig:check-missing-boundary}
\end{figure}

% ============================================================
\section{查找并修复天线违例}
\subsection*{Finding and Fixing Antenna Violations}

芯片制造中，栅氧易被静电放电损伤。多层金属化过程中导线累积的静电荷可能损坏器件或导致芯片完全失效。静电电荷泄放入器件的现象称为天线（antenna）或电荷收集天线问题。

为防止天线问题，工具验证每个输入 pin 的金属天线面积除以栅极面积小于晶圆厂给出的最大天线比：
\[
\frac{\text{antenna-area}}{\text{gate-area}} < \text{max-antenna-ratio}
\]
该检查基于 block 的全局与层相关天线规则，以及单元的天线属性。

天线流程步骤：
\begin{enumerate}
  \item 定义天线规则；
  \item 指定 pin 与 port 的天线属性；
  \item 分析并修复天线违例。
\end{enumerate}

\subsection{定义天线规则}
\subsubsection*{Defining Antenna Rules}

天线规则定义如何计算 block 中网络的最大天线比，以及 pin 的天线比。必须用 \cmd{define\_antenna\_rule} 定义全局天线规则；也可用 \cmd{define\_antenna\_layer\_rule} 指定层相关规则。无层规则时使用全局规则。

相关主题：
\begin{itemize}
  \item 计算最大天线比
  \item 计算 Pin 的天线比
\end{itemize}

\subsubsection{计算最大天线比}
\subsubsection*{Calculating the Maximum Antenna Ratio}

须指定：
\begin{itemize}
  \item 二极管提供多少保护（diode protection mode）——\cmd{define\_antenna\_rule} 的 \opt{-diode\_mode}（必选），见「设置 Diode Protection Mode」；
  \item 在二极管保护下如何计算天线比（diode ratio vector）——\cmd{define\_antenna\_layer\_rule} 的 \opt{-diode\_ratio}（必选），见「指定 Diode Ratio Vector」。
\end{itemize}

\subsubsection{设置 Diode Protection Mode}
\subsubsection*{Setting the Diode Protection Mode}

Diode protection mode 规定二极管提供的保护程度。先进工艺设计常使用不同 gate oxide class 表示多种栅氧厚度。工具最多支持四种：gate class 0、1、2、3。无 gate class 数据的 pin 视为 class 0。对使用多种栅氧厚度的设计，须使用支持多厚度的模式：14、16、18。

用 \cmd{define\_antenna\_rule -diode\_mode} 指定（必选）。表~\ref{tab:diode-protection-mode} 定义各模式如何影响网络最大天线比；各二极管最大天线比的计算公式见表~\ref{tab:antenna-ratio-calc}，其中使用 \opt{-diode\_ratio} 向量值。

\begin{table}[htbp]
\centering
\caption{Diode Protection Mode 设置（原书 Table 28）}
\label{tab:diode-protection-mode}
\begin{tabular}{@{}cp{0.72\textwidth}@{}}
\toprule
模式 & 定义 \\
\midrule
0 & 二极管不提供任何保护。 \\
1 & 二极管提供无限保护。此模式下最高金属层不检查天线违例，因为最高层形成时输入--输出连接已完成。 \\
2 & 保护有限；多二极管时取各二极管最大天线比的最大值。 \\
3 & 保护有限；多二极管时取各二极管最大天线比之和。 \\
4 & 保护有限；多二极管时用各二极管 diode-protection 值之和计算最大天线比。 \\
5 & 保护有限；用所有二极管的最大 diode-protection 值计算等效栅极面积。 \\
6 & 保护有限；用所有二极管 diode-protection 值之和计算等效栅极面积。 \\
7 & 保护有限；用最大 diode-protection 值计算等效金属面积。 \\
8 & 保护有限；用 diode-protection 值之和计算等效金属面积。 \\
9 & 保护有限；按最大二极管保护做缩放。 \\
10 & 保护有限；按总二极管保护做缩放。 \\
11 & 保护有限；按最大二极管保护做缩放。 \\
12 & 保护有限；按总二极管保护做缩放。 \\
13 & 已废弃。 \\
14 & 保护有限；两套独立计算公式。支持 gate class。 \\
15 & 仅内部使用。 \\
16 & 保护有限；多二极管时用 diode-protection 值之和计算最大天线比。支持 gate class。 \\
17 & 保护有限；用 diode-protection 值之和计算等效栅极面积；输入与输出 pin 分别计算。 \\
18 & 保护有限；多二极管时用 diode-protection 值之和计算最大天线比。支持 gate class，且计算取决于 pin 是否连二极管。 \\
\bottomrule
\end{tabular}
\end{table}

\subsubsection{指定 Diode Ratio Vector}
\subsubsection*{Specifying the Diode Ratio Vector}

用 \cmd{define\_antenna\_layer\_rule -diode\_ratio} 指定用于计算最大天线比的值。未指定时工具使用 \texttt{\{0 0 1 0 0\}}。实际用法取决于 \cmd{define\_antenna\_rule -diode\_mode}。

表~\ref{tab:diode-ratio-format} 给出各模式的向量格式；表~\ref{tab:antenna-ratio-calc} 说明如何用向量值计算。其中：
\begin{itemize}
  \item \texttt{dp}：输出 pin 的 diode protection 值（见「指定天线属性」）；
  \item \texttt{layerMaxRatio}：该层最大天线比，由 \cmd{define\_antenna\_layer\_rule -ratio} 或全局规则的 \opt{-metal\_ratio}/\opt{-cut\_ratio} 指定。
\end{itemize}

示例见后文「模式 2/3/4」「模式 5/6」「模式 7/8」「模式 14 多栅氧厚度」各小节。

\begin{table}[htbp]
\centering
\caption{Diode Ratio Vector 格式（原书 Table 29）}
\label{tab:diode-ratio-format}
\begin{tabular}{@{}cp{0.72\textwidth}@{}}
\toprule
模式 & 格式 \\
\midrule
2--10, 17 & \texttt{\{v0 v1 v2 v3 [v4]\}}；\texttt{v4} 为可选保护上限，省略视为 0（无上限）。 \\
11, 12 & \texttt{\{\{v0 v1 v2 v3 [v4]\}\{s0 s1 s2 s3 s4 s5\}\}}；第二向量为缩放值。 \\
14 & 每个 gate class 一个 10 值向量（共 4 个）。 \\
16 & 每个 gate class 一个 8 值向量（共 4 个）。 \\
18 & 每个 gate class 一个 12 值向量（共 4 个）。 \\
\bottomrule
\end{tabular}
\end{table}

\begin{table}[htbp]
\centering
\caption{基于 Diode Mode 的天线比计算（原书 Table 30）}
\label{tab:antenna-ratio-calc}
\small
\begin{tabularx}{\textwidth}{@{}c >{\raggedright\arraybackslash}X@{}}
\toprule
模式 & 计算 \\
\midrule
0, 1 & 不使用 \\
2, 3, 4 &
若 \texttt{dp>v0} 且 \texttt{v4$\neq$0}：\texttt{antenna\_ratio = min(((dp+v1)*v2+v3), v4)}；\newline
若 \texttt{dp>v0} 且 \texttt{v4=0}：\texttt{(dp+v1)*v2+v3}；\newline
若 \texttt{dp$\leq$v0}：\texttt{layerMaxRatio} \\
5 & \texttt{antenna\_ratio = metal\_area/(gate\_area+equi\_gate\_area)}；\newline
若 \texttt{max\_diode\_protection>v0} 且 \texttt{v4$\neq$0}：
\texttt{equi\_gate\_area = min(((max\_diode\_protection+v1)*v2+v3), v4)}；\newline
若 \texttt{v4=0}：\texttt{(max\_diode\_protection+v1)*v2+v3}；\newline
若 \texttt{max\_diode\_protection$\leq$v0}：\texttt{equi\_gate\_area=0} \\
6 & 同 5，但用 \texttt{total\_diode\_protection} \\
7 & \texttt{antenna\_ratio = (metal\_area-equi\_metal\_area)/gate\_area}；\newline
\texttt{equi\_metal\_area} 由 \texttt{max\_diode\_protection} 按与模式 5 相同的 v0--v4 规则计算；\newline
若 \texttt{equi\_metal\_area>metal\_area} 则取 \texttt{metal\_area} \\
8 & 同 7，但用 \texttt{total\_diode\_protection} \\
9 & \texttt{antenna\_ratio = scale*metal\_area/gate\_area}；\newline
若 \texttt{max\_diode\_protection>v0}：
\texttt{scale = max(1/((max\_diode\_protection+v1)*v2+v3), v4)}；\newline
否则 \texttt{scale=1.0} \\
10 & 同 9，但用 \texttt{total\_diode\_protection} \\
11 & \texttt{antenna\_ratio = scale*metal\_area/gate\_area}；\newline
按 \texttt{max\_diode\_protection} 相对 v0--v4 的区间选取 \texttt{s0}--\texttt{s5} \\
12 & 同 11，但用 \texttt{total\_diode\_protection} \\
13 & N/A \\
14 & 连二极管：两式同时成立\newline
\texttt{(metal/via)/gate $\leq$ (dp*v1+v2)} 且
\texttt{(metal/via)/(v5*gate+v6*dp) $\leq$ v7}；\newline
不连二极管：\texttt{(metal/via)/gate $\leq$ v3} 且
\texttt{(metal/via)/(v8*gate) $\leq$ v9} \\
15 & N/A（内部使用） \\
16 & 允许最大天线比：同模式 2/3/4 的 dp/v0--v4 规则；\newline
若 \texttt{gate-area $\leq$ v5}：\texttt{antenna-ratio = antenna-area/(gate-area+v6*dp)}；\newline
否则 \texttt{antenna-area/(gate-area+v6*dp+v7)} \\
17 & 同模式 6（等效栅极面积），输入/输出 pin 分别计算 \\
18 & 连二极管：同模式 16；不连二极管：按 gate-area 与 v8--v11 取允许最大比，\newline
\texttt{antenna-ratio = antenna-area/gate-area} \\
\bottomrule
\end{tabularx}
\end{table}

\paragraph{模式 2、3、4 示例}
\subsubsection*{Example for Diode Modes 2, 3, and 4}

假设天线面积按单层表面面积（antenna mode 1）计算，M1 的 diode ratio vector 为 \texttt{\{0.7 0.0 200 2000\}}，\texttt{layerMaxRatio}=400，同一 M1 网连接二极管 A（保护值 0.5）、B（1.0）、C（1.5）；未指定 \texttt{v4} 时按 0。

\begin{lstlisting}
fc_shell> define_antenna_rule -mode 1 -diode_mode diode_mode \
   -metal_ratio 400 -cut_ratio 20
fc_shell> define_antenna_layer_rule -mode 1 -layer "M1" \
   -ratio 400 -diode_ratio {0.7 0.0 200 2000}
\end{lstlisting}

各二极管最大天线比按表~\ref{tab:antenna-ratio-calc} 与表~\ref{tab:max-ratio-per-diode} 计算。

\begin{table}[htbp]
\centering
\caption{各二极管最大天线比计算（原书 Table 31）}
\label{tab:max-ratio-per-diode}
\begin{tabular}{@{}clp{0.55\textwidth}@{}}
\toprule
二极管 & 保护值 & 最大天线比 \\
\midrule
A & 0.5 & 400（保护值 $<$ v0=0.7，用 \texttt{layerMaxRatio}） \\
B & 1.0 & 2200（$(1.0+0.0)*200+2000$） \\
C & 1.5 & 2300（$(1.5+0.0)*200+2000$） \\
\bottomrule
\end{tabular}
\end{table}

网络最大天线比：
\begin{itemize}
  \item 模式 2：取各二极管最大值 $=2300$；
  \item 模式 3：求和 $=400+2200+2300=4900$；
  \item 模式 4：用保护值之和代入公式，\((0.5+1.0+1.5)*200+2000=2600\)。
\end{itemize}

\paragraph{模式 5、6 示例}
\subsubsection*{Example for Diode Modes 5 and 6}

假设天线面积按单层模式表面面积（antenna mode 1）计算，M1 的 diode ratio vector 为 \texttt{\{0 0 1 0\}}（默认），M1 的 \texttt{layerMaxRatio} 为 400，同一 M1 网络连接二极管 A（保护值 0.5）、B（1.0）、C（1.5）；未指定 \texttt{v4} 时按 0 处理。

\begin{lstlisting}
fc_shell> define_antenna_rule -mode 1 -diode_mode diode_mode \
   -metal_ratio 400 -cut_ratio 20
\end{lstlisting}

因使用全局 \texttt{layerMaxRatio} 与默认向量，无需层相关规则。模式 5/6 的最大天线比为
\texttt{metal\_area / (gate\_area + equi\_gate\_area)}，其中
\texttt{equi\_gate\_area = (diode\_protection + v1) * (v2 + v3)}。
\begin{itemize}
  \item 模式 5：用最大保护值 1.5，得 \texttt{metal\_area/(gate\_area + 1.5)}；
  \item 模式 6：用总保护值 \(0.5+1.0+1.5=3.0\)，得 \texttt{metal\_area/(gate\_area + 3.0)}。
\end{itemize}

\paragraph{模式 7、8 示例}
\subsubsection*{Example for Diode Modes 7 and 8}

假设单层表面面积模式，M1 向量为 \texttt{\{0.7 0.0 150 800\}}，\texttt{layerMaxRatio}=400，二极管保护值同前。

\begin{lstlisting}
fc_shell> define_antenna_rule -mode 1 -diode_mode diode_mode \
   -metal_ratio 400 -cut_ratio 20
fc_shell> define_antenna_layer_rule -mode 1 -layer "M1" \
   -ratio 400 -diode_ratio {0.7 0.0 150 800}
\end{lstlisting}

模式 7/8 计算为 \texttt{(metal\_area - equi\_metal\_area) / gate\_area}：
\begin{itemize}
  \item 模式 7：最大保护 1.5，得 \texttt{(metal\_area - 1025) / gate\_area}；
  \item 模式 8：总保护 3.0，得 \texttt{(metal\_area - 1250) / gate\_area}。
\end{itemize}

\paragraph{模式 14（多栅氧厚度）示例}
\subsubsection*{Example for Diode Mode 14 With Multiple Gate Oxide Thicknesses}

假设单层表面面积模式，设计含两种栅氧厚度。gate class 0 的 M1 向量为
\texttt{\{1 0 1e9 1e9 0 1 2 285 1 285\}}，gate class 1 为
\texttt{\{1 0 1e9 1e9 1 0 2 165 1 28\}}，\texttt{layerMaxRatio}=285。
模式 14 每个 gate class 需 10 个值；即使仅用两类，定义时仍须给出全部 40 个值。未用 class 填很大比率以禁用检查。

\begin{lstlisting}
fc_shell> define_antenna_rule -mode 1 -diode_mode 14 \
   -metal_ratio 285 -cut_ratio 10
fc_shell> define_antenna_layer_rule -mode 1 -layer "M1" \
   -ratio 285 -diode_ratio {1 0 1e9 1e9 0 1 2 285 1 285 \
                            1 0 1e9 1e9 1 0 2 165 1 28 \
                            1 0 1e9 1e9 0 1 2 1e9 1 1e9 \
                            1 0 1e9 1e9 0 1 2 1e9 1 1e9 }
\end{lstlisting}

向量各分量含义：
\begin{itemize}
  \item V0：公式 1 使用的栅极面积范围（0=当前 class；1=全部 class 总和）；
  \item V1：连二极管时公式 1 的 diode protection 乘数；
  \item V2：连二极管时公式 1 的偏移；
  \item V3：不连二极管时公式 1 的比率要求；
  \item V4：公式 2 的栅极面积范围（同 V0）；
  \item V5--V7：连二极管时公式 2 的栅极乘数、保护乘数与比率要求；
  \item V8--V9：不连二极管时公式 2 的栅极乘数与天线比要求。
\end{itemize}

模式 14 公式（表~\ref{tab:antenna-ratio-calc} 续）：若连二极管，须同时满足
\texttt{(metal/via area)/gate-area <= (dp*v1+v2)} 与
\texttt{(metal/via area)/(v5*gate-area+v6*dp) <= v7}；
若不连二极管，须同时满足
\texttt{(metal/via area)/gate-area <= v3} 与
\texttt{(metal/via area)/(v8*gate-area) <= v9}。

\subsubsection{计算 Pin 的天线比}
\subsubsection*{Calculating the Antenna Ratio for a Pin}

工具支持多种天线比（antenna-area/gate-area）计算方式。通过指定下列信息控制：
\begin{itemize}
  \item 如何计算天线面积（antenna area mode）；
  \item 计算时考虑哪些金属段（antenna recognition mode）。
\end{itemize}
用 \cmd{define\_antenna\_rule} 与 \cmd{define\_antenna\_layer\_rule} 的 \opt{-mode} 指定（二者的必选选项），见「设置天线模式」。

\subsection{设置天线模式}
\subsubsection*{Setting the Antenna Mode}

天线模式通过以下两方面控制天线比计算：
\begin{itemize}
  \item \textbf{天线面积计算方式（antenna area mode）}
  \begin{itemize}
    \item 表面面积：\(W \times L\)；
    \item 侧壁面积：\((W+L)\times 2\times\text{thickness}\)。
\begin{noteBox}
使用侧壁面积时，须在 technology file 各 Layer 段中用 \texttt{unitMinThickness}、\texttt{unitNomThickness}、\texttt{unitMaxThickness} 定义金属厚度。
\end{noteBox}
  \end{itemize}
  \item \textbf{天线识别模式（antenna recognition mode）}
  \begin{itemize}
    \item \textbf{Single-layer mode}：仅考虑当前层金属段，忽略下层。可布线性最好。\\
          \texttt{antenna\_ratio = 该层相连金属面积 / 总栅极面积}
    \item \textbf{Accumulated-ratio mode}：考虑当前层及与输入 pin 相连的下层段。\\
          \texttt{antenna\_ratio = 当前层与下层各 single-layer 比率之和}
    \item \textbf{Accumulated-area mode}：考虑当前层与下层金属面积。\\
          \texttt{antenna\_ratio = 当前及下层相连金属面积 / 总栅极面积}
  \end{itemize}
\end{itemize}

用 \opt{-mode} 指定天线模式（必选）。表~\ref{tab:antenna-mode} 列出取值。

\begin{table}[htbp]
\centering
\caption{天线模式设置（原书 Table 32）}
\label{tab:antenna-mode}
\begin{tabular}{@{}ccc@{}}
\toprule
天线模式值 & 面积计算方式 & 天线识别模式 \\
\midrule
1 & Surface & Single-layer \\
2 & Surface & Accumulated-ratio \\
3 & Surface & Accumulated-area \\
4 & Sidewall & Single-layer \\
5 & Sidewall & Accumulated-ratio \\
6 & Sidewall & Accumulated-area \\
\bottomrule
\end{tabular}
\end{table}

\subsubsection{天线识别模式示例}
\subsubsection*{Antenna Recognition Mode Example}

图~\ref{fig:antenna-recog-layout} 给出含侧视的布局示例；表~\ref{tab:antenna-recog-ratios} 列出各识别模式下的天线比。

\begin{figure}[htbp]
\centering
\fbox{\parbox{0.85\textwidth}{\centering\vspace{1.5cm}
\textit{（原书 Figure 89：Layout Example for Antenna Recognition Modes）}\\[0.3em]
\small 请参阅原版 PDF 插图。
\vspace{1.5cm}}}
\caption{天线识别模式布局示例}
\label{fig:antenna-recog-layout}
\end{figure}

\begin{table}[htbp]
\centering
\caption{天线识别模式与比率（原书 Table 33）}
\label{tab:antenna-recog-ratios}
\small
\begin{tabularx}{\textwidth}{@{}l >{\raggedright\arraybackslash}X@{}}
\toprule
考虑的段 & 天线比 \\
\midrule
Single-layer &
M1：Gate1$=$M1b/Gate1；Gate2$=$M1c/Gate2；Gate3$=$M1d/Gate3。\newline
M2：Gate1$=$M2b/Gate1；Gate2$=$M2d/Gate2；Gate3$=$M2e/Gate3。\newline
M3：Gate1,2$=$(M3b+M3c)/(Gate1+Gate2)；Gate3$=$M3d/Gate3。\newline
M4：Gate1,2,3$=$(M4a+M4b)/(Gate1+Gate2+Gate3)。 \\
\midrule
Accumulated-ratio &
M1 同 single-layer。\newline
M2：Gate1$=$M1b/Gate1$+$M2b/Gate1（Gate2/3 类似累加）。\newline
M3/M4：在下层比率上继续累加当前层 single-layer 比率。 \\
\midrule
Accumulated-area &
M1 同 single-layer。\newline
M2：Gate1$=$(M1b+M2b)/Gate1 等。\newline
M3：Gate1,2 为多层金属面积之和 / (Gate1+Gate2)；Gate3 类似。\newline
M4：全部金属面积 / (Gate1+Gate2+Gate3)。 \\
\bottomrule
\end{tabularx}
\end{table}

\subsection{指定天线属性}
\subsubsection*{Specifying Antenna Properties}

标准单元与 hard macro 的天线属性一般在参考库的 frame view 中定义。可为未在库中定义属性的单元设置默认值：
\begin{itemize}
  \item \appopt{route.detail.default\_diode\_protection}：标准单元输出 pin 缺省 diode protection；
  \item \appopt{route.detail.default\_gate\_size}：标准单元输入 pin 缺省 gate size；
  \item \appopt{route.detail.default\_port\_external\_antenna\_area}：顶层 port 缺省天线面积；
  \item \appopt{route.detail.default\_port\_external\_gate\_size}：顶层 port 缺省 gate size；
  \item \appopt{route.detail.macro\_pin\_antenna\_mode}：macro pin 的天线处理方式；
  \item \appopt{route.detail.port\_antenna\_mode}：顶层 port 的天线处理方式。
\end{itemize}

层次化流程中若为 block 创建 abstraction，须用 \cmd{derive\_hier\_antenna\_property} 从 block 提取天线信息供上一层次使用。

\begin{noteBox}
若 macro 并非所有层都定义了层次化天线属性，Zroute 视为数据不完整，跳过天线分析并给出 ZRT-311 警告。参见 SolvNet 文章 027178。
\end{noteBox}

\begin{seeAlsoBox}
在 hard macro 上标注天线属性（Annotating Antenna Properties on Hard Macro Cells）。
\end{seeAlsoBox}

\subsection{分析并修复天线违例}
\subsubsection*{Analyzing and Fixing Antenna Violations}

若设计库含天线规则，Zroute 在详细布线期间自动分析并修复天线违例。要禁用，将 \appopt{route.detail.antenna} 设为 \cmd{false}。天线规则与其他设计规则一样在详细布线中并发检查与修正，以减少迭代、缩短总运行时间。

默认行为：
\begin{itemize}
  \item 对全部时钟与信号网检查并修复；可用 \appopt{route.detail.skip\_antenna\_fixing\_for\_nets} 禁止对特定网修复（仍分析并报告）；
  \item 不对电源/地网检查或修复；将 \appopt{route.detail.check\_antenna\_on\_pg} 设为 \cmd{true} 可启用；
  \item 在第二次迭代（初始布线完成且基本 DRC 已修复后）开始修复；可用 \appopt{route.detail.antenna\_on\_iteration} 更改；
  \item 通过 \textbf{layer hopping} 修复：把大金属多边形拆到更高层，或（对最高层有限保护情形）下移到更低层。
\end{itemize}

Layer hopping 类型：
\begin{itemize}
  \item 用更高层金属段打断天线：对 metal-N 违例，在靠近栅极处插入小段 metal-(N+1)，降低剩余 metal-N 比率。不适用于最高层且输出保护有限的情况。
  \item 下移到更低层：主要用于最高层违例；可能增加 RC 与时序延迟。
\end{itemize}

也可插入二极管修复。将 \appopt{route.detail.insert\_diodes\_during\_routing} 设为 \cmd{true} 启用。要强制仅用二极管，将 \appopt{route.detail.hop\_layers\_to\_fix\_antenna} 设为 \cmd{false}。若两者皆启用，默认先尝试 layer hopping；要将偏好改为二极管，设 \appopt{route.detail.antenna\_fixing\_preference} 为 \texttt{use\_diodes}。

删除修复天线所不需要的冗余二极管：\appopt{route.detail.delete\_redundant\_diodes\_during\_routing}=\cmd{true}。Spare 二极管不删除；详细布线插入的 port 保护二极管也不删除。

每轮详细布线结束会报告天线违例，例如：
\begin{lstlisting}
DRC-SUMMARY:
   @@@@@@@ TOTAL VIOLATIONS =      506
    @@@@ Total number of instance ports with antenna violations = 1107
\end{lstlisting}

检查天线违例：
\begin{lstlisting}
fc_shell> check_routes -antenna true
\end{lstlisting}

\subsection{详细布线期间插入二极管}
\subsubsection*{Inserting Diodes During Detail Routing}

通过在靠近被保护栅极的网络上插入反向偏置二极管，为累积电荷提供放电通路，且正常工作时不应导通电流。

启用：\appopt{route.detail.insert\_diodes\_during\_routing}=\cmd{true}。

控制选项：
\begin{itemize}
  \item \appopt{route.detail.antenna\_fixing\_preference}=\texttt{use\_diodes}：优先二极管而非 hopping；
  \item \appopt{route.detail.hop\_layers\_to\_fix\_antenna}=\cmd{false}：强制二极管修复；
  \item 默认可新增二极管或使用 spare；用 \appopt{route.detail.diode\_preference} 设为 \texttt{new}/\texttt{spare}/\texttt{none}；
  \item \appopt{route.detail.diode\_insertion\_mode}=\texttt{new}/\texttt{spare}/\texttt{new\_and\_spare}：强制方法；
\begin{noteBox}
要利用 spare 二极管，须在标准单元布局前/后、天线违例可能出现区域布线之前加入 spare。
\end{noteBox}
  \item \appopt{route.detail.diode\_libcell\_names}：限制可用二极管（仅写单元名，勿含库名）；
  \item \appopt{route.detail.reuse\_filler\_locations\_for\_diodes}=\cmd{false}：禁止复用 filler 位置。
\end{itemize}

层次与 voltage area：
\begin{itemize}
  \item Pin 违例：二极管插在与违例 pin 相同逻辑层次；
  \item 顶层 port 违例：默认顶层；若 port 属 voltage area，可将 \appopt{route.detail.use\_lower\_hierarchy\_for\_port\_diodes} 设为 \cmd{true}。
\end{itemize}

\subsubsection{详细布线后插入二极管}
\subsubsection*{Inserting Diodes After Detail Routing}

详细布线后可用 \cmd{create\_diodes} 显式指定违例与约束。根据约束与 \appopt{route.detail.diode\_insertion\_mode}，插入新二极管或复用 spare。

须用列表指定：port、单元实例、二极管参考单元、数量、连接最高层、相对 pin 的最大距离。距离内无法插入/复用则跳过。

\begin{lstlisting}
{port_name instance_name diode_reference number_of_diodes
   max_routing_layer max_routing_distance}
\end{lstlisting}

\begin{noteBox}
将单元实例指定为顶层 block 名，可为顶层 port 插入二极管。
\end{noteBox}

相关应用选项：\appopt{route.detail.diode\_insertion\_mode}、\appopt{route.detail.reuse\_filler\_locations\_for\_diodes}、\appopt{route.detail.use\_lower\_hierarchy\_for\_port\_diodes}。

示例（CI 的 A port，2 个 Adiode，最高 M5，2.5\,µm）：
\begin{lstlisting}
fc_shell> create_diodes {{A CI Adiode 2 M5 2.5}}
\end{lstlisting}

% ============================================================
\section{插入冗余 Via}
\subsection*{Inserting Redundant Vias}

冗余 via 插入是 Zroute 在整个布线流程中支持的重要 DFM 功能。各布线阶段并发优化 via 数量与线长。结果以冗余 via 转换率（单 via 转为冗余 via 的百分比）衡量；亦应关注未优化单 via 数量——越少通常对 DFM 越有利。

相关主题：
\begin{itemize}
  \item 在时钟网上插入冗余 Via
  \item 在信号网上插入冗余 Via
\end{itemize}

\subsection{在时钟网上插入冗余 Via}
\subsubsection*{Inserting Redundant Vias on Clock Nets}

Zroute 可在时钟布线期间或之后插入冗余 via。布线后插入见「布线后冗余 Via 插入」。

时钟布线期间步骤：
\begin{enumerate}
  \item 用 \cmd{create\_routing\_rule} 在非默认布线规则中指定冗余 via（见「指定非默认 Via」）。选择时须兼顾 DFM 与可布线性。NDR 还可定义更严线宽/间距与 tapering 距离。
  \item 用 \cmd{set\_clock\_routing\_rules} 将 NDR 赋给时钟网：
\begin{lstlisting}
fc_shell> set_clock_routing_rules -rules clock_via_rule
\end{lstlisting}
  \item 运行 CTS：
\begin{lstlisting}
fc_shell> synthesize_clock_trees
\end{lstlisting}
  \item 布线时钟网：
\begin{lstlisting}
fc_shell> route_group -all_clock_nets \
   -reuse_existing_global_route true
\end{lstlisting}
        全局布线预留冗余 via 空间，详细布线插入冗余 via。
\end{enumerate}

\subsection{在信号网上插入冗余 Via}
\subsubsection*{Inserting Redundant Vias on Signal Nets}

可采用：
\begin{itemize}
  \item 布线后冗余 via 插入；
  \item 并发 soft-rule 冗余 via 插入；
  \item 接近 100\% 冗余 via 插入。
\end{itemize}

一般先做布线后插入；若转换率 $\ge$80\% 可再尝试并发 soft-rule；若 $\ge$90\% 可再尝试接近 100\% 插入。

\begin{noteBox}
转换率升高会使设计规则更难收敛，并可能降低信号完整性；接近 100\% 仅用于确需极高转换率的 block，且可能需调整 floorplan utilization 以留出空间。
\end{noteBox}

相关主题：查看默认 Via 映射表；定义定制 Via 映射表；布线后插入；并发 soft-rule；接近 100\%；时序保持；报告转换率。

\subsubsection{查看默认 Via 映射表}
\subsubsection*{Viewing the Default Via Mapping Table}

默认 Zroute 从 technology file 读取默认 contact code 并生成优化映射表。查看：
\begin{lstlisting}
fc_shell> add_redundant_vias -list_only true
\end{lstlisting}
多数情况下定制映射表优于默认表。若已用 \cmd{add\_via\_mapping} 创建定制表，该命令显示定制表。

示例（原书 Example 25）：
\begin{lstlisting}
fc_shell> add_redundant_vias -list_only
...
Redundant via optimization will attempt to replace the following vias:
    VIA12SQ_C    -> VIA12SQ_C_2x1    VIA12SQ_C_2x1(r) VIA12SQ_C_1x2 ...
    VIA12SQ_C(r) -> VIA12SQ_C_2x1    VIA12SQ_C_2x1(r) VIA12SQ_C_1x2 ...
...
\end{lstlisting}

\subsubsection{定义定制 Via 映射表}
\subsubsection*{Defining a Customized Via Mapping Table}

使用 \cmd{add\_via\_mapping}。至少用 \opt{-from}/\opt{-to} 指定源 via 与替换 via（须为 technology file 中的 via 或 \cmd{create\_via\_def} 创建的设计专用 via）。\opt{-from} 须为 simple via 或 via array；\opt{-to} 可为 simple、simple array 或 custom via。

细化选项：
\begin{itemize}
  \item \opt{-weight}：优先级 1--10（默认 1）；高权重优先使用；
  \item \opt{-transform}：\texttt{all}（默认可旋转/翻转）、\texttt{rotate}、\texttt{flip}、\texttt{none}。
\end{itemize}

映射保存在设计库 via 映射表中。覆盖已有映射须 \opt{-force}。

\begin{lstlisting}
fc_shell> add_via_mapping \
   -from {VIA12 1x1} -to {VIA12T 2x1} -weight 5
fc_shell> add_via_mapping \
   -from {VIA12 1x1} -to {VIA12 2x1} -weight 1
\end{lstlisting}

用 \cmd{report\_via\_mapping} 查看（\opt{-from}/\opt{-to} 可过滤）。用 \cmd{remove\_via\_mappings} 移除（\opt{-all}/\opt{-from}/\opt{-to}）。

\subsubsection{使用 Via 映射表子集}
\subsubsection*{Using a Subset of the Via Mapping Table for Redundant Via Insertion}

通过应用选项按层选择权重组：
\begin{itemize}
  \item \appopt{route.common.redundant\_via\_include\_weight\_group\_by\_layer\_name}：仅用指定权重组；未列出的层不做插入；
  \item \appopt{route.common.redundant\_via\_exclude\_weight\_group\_by\_layer\_name}：排除指定权重组；
  \item 两者同时设置时，在 include 集合中再排除 exclude 集合。
\end{itemize}

\subsubsection{布线后冗余 Via 插入}
\subsubsection*{Postroute Redundant Via Insertion}

使用 \cmd{add\_redundant\_vias}：可将单 cut 换成多 cut 阵列、换成其他 contact code 的单 cut，或换成不同多 cut 阵列。插入时详细布线检查邻近 partition 的设计规则以减少违例。

默认对所有网插入；\opt{-nets} 可限制网络。完成后重新检查并修复 DRC。

若转换率不足，可用 \opt{-effort high}（约再提升 3--5\%，通过移动 via 腾空间）；在 45\,nm 及以下可能不利于光刻友好图案，此时宜改用并发 soft-rule。也可将 \appopt{route.detail.optimize\_wire\_via\_effort\_level} 设为 \texttt{high} 以减少单 via、缩短线长。

初始布线后插入完成后，设 \appopt{route.common.post\_detail\_route\_redundant\_via\_insertion}，以便后续详细/ECO 布线后自动插入，维持转换率。

\subsubsection{并发 Soft-Rule 冗余 Via 插入}
\subsubsection*{Concurrent Soft-Rule-Based Redundant Via Insertion}

在布线期间为冗余 via 预留空间以提高转换率；可用于初始与 ECO 布线。实际插入仍须事后运行 \cmd{add\_redundant\_vias}。

\begin{noteBox}
预留空间会增加布线运行时间；仅当布线后方法转换率已 $\ge$80\% 且需进一步提高时使用。
\end{noteBox}

步骤：
\begin{enumerate}
  \item （可选）定义 via 映射表；
  \item 启用并发 soft-rule：
\begin{lstlisting}
fc_shell> set_app_options \
   -name route.common.concurrent_redundant_via_mode \
   -value reserve_space
\end{lstlisting}
        用 \appopt{route.common.concurrent\_redundant\_via\_effort\_level} 控制努力（默认 low；high 也影响详细布线，可能影响 DRC 收敛）。若在 \cmd{place\_opt} 前启用，拥塞估计会考虑冗余 via。\\
        ECO：设 \appopt{route.common.eco\_route\_concurrent\_redundant\_via\_mode}=\texttt{reserve\_space}，并用对应 effort 选项。
\begin{noteBox}
ECO 期间使用可能影响时序与 DRC 收敛；一般仅当初始布线用了接近 100\% 插入时才在 ECO 用 soft-rule 维持转换率。
\end{noteBox}
  \item 布线 block（预留空间并修复硬规则）；
  \item 运行布线后 \cmd{add\_redundant\_vias}。
\end{enumerate}

\subsubsection{接近 100\% 冗余 Via 插入}
\subsubsection*{Near 100 Percent Redundant Via Insertion}

将冗余 via 作为硬设计规则处理。仅支持初始布线，不支持 ECO；初始用此法后，ECO 应改用 soft-rule 以保持转换率。

\begin{noteBox}
对拥塞 block 可能大幅增加运行时间；仅当 soft-rule 已达 $\ge$90\% 且仍需提高时使用。
\end{noteBox}

\begin{lstlisting}
fc_shell> set_app_options \
   -name route.common.concurrent_redundant_via_mode \
   -value insert_at_high_cost
\end{lstlisting}

布线期间插入冗余 via 并修复硬规则；优先级通常与其他硬规则相同，若详细布线 DRC 不收敛会自动放宽冗余 via 约束。

\subsubsection{冗余 Via 插入时保持时序}
\subsubsection*{Preserving Timing During Redundant Via Insertion}

插入冗余 via 会改变时序：短网因电容增加变慢，长网因电阻下降变快。时序保持模式通过对关键网禁止插入来避免影响时序。

在 \cmd{add\_redundant\_vias} 上使用：
\begin{itemize}
  \item \opt{-timing\_preserve\_setup\_slack\_threshold}
  \item \opt{-timing\_preserve\_hold\_slack\_threshold}
  \item \opt{-timing\_preserve\_nets}
\end{itemize}

应在流程末尾、常规插入已收敛转换率与时序 QoR 之后使用；过早使用可能因关键网过多而严重降低转换率。

\subsection{报告冗余 Via 率}
\subsubsection*{Reporting Redundant Via Rates}

插入后 Zroute 生成报告，包含：
\begin{itemize}
  \item 非默认 via 的总优化转换率；
  \item 各层优化转换率（含 double via 与 DFM-friendly bar via），并分别相对总 via 数与可优化 routed via 数给出；
\begin{noteBox}
使用 bar via 时优化转换率参考价值有限。
\end{noteBox}
  \item 各层按 weight 的优化 via 分布（高于某权重的转换率须对各权重求和）；无权重记为 ``Un-optimized''；
  \item block 总 double via 转换率。
\end{itemize}

\begin{noteBox}
\cmd{report\_design} 的冗余 via 率与 Zroute/\cmd{check\_routes} 不同：前者含 PG 与信号 via，后者仅信号 via，且基于冗余 via 映射。
\end{noteBox}

示例（原书 Example 27）：
\begin{lstlisting}
Total optimized via conversion rate = 96.94% (1401030 / 1445268 vias)
Layer V01        = 41.89% (490617 / 1171301 vias)
    Weight 10    = 9.64% (112869 vias)
    Weight 5     = 32.25% (377689 vias)
    ...
Total double via conversion rate = 46.69% (2158006 / 4622189 vias)
\end{lstlisting}

% ============================================================
\section{优化线长与 Via 数量}
\subsection*{Optimizing Wire Length and Via Count}

详细布线时，Zroute 在存在 DRC 违例的区域优化线长与 via 数量，但不优化无违例区域。

为提高良率，在详细布线与冗余 via 插入之后，可用 \cmd{optimize\_routes} 独立优化线长与 via 数量。

默认按总体代价选择重布网络；对所选网决定重布全部形状或部分形状。用 \opt{-nets} 指定网络；可用 \opt{-reroute\_all\_shapes\_in\_nets} 强制重布该网全部形状。

默认最多 40 次详细布线迭代以修复优化后违例；可用 \opt{-max\_detail\_route\_iterations} 控制。

% ============================================================
\section{减小临界面积}
\subsection*{Reducing Critical Areas}

临界面积是这样一块区域：若随机颗粒缺陷中心落在该处，会导致电路失效、降低良率。导电缺陷造成短路，非导电缺陷造成开路。

相关主题：
\begin{itemize}
  \item 通过 wire spreading 减小短路临界面积；
  \item 通过 wire widening 减小开路临界面积。
\end{itemize}

\subsection{执行 Wire Spreading}
\subsubsection*{Performing Wire Spreading}

详细布线与冗余 via 插入后，可用 \cmd{spread\_wires} 增大平均线间距，减小短路临界面积。默认对同层信号线：
\begin{itemize}
  \item 在优选方向疏导半个 pitch（用 \opt{-pitch} 以层 pitch 倍数修改，如 \texttt{-pitch 1.5}）；
  \item 最小 jog 长度为两倍层 pitch（\opt{-min\_jog\_length}，层 pitch 整数倍）；
  \item 最小 jog 间距为最小层间距加半个层 pitch（\opt{-min\_jog\_spacing\_by\_layer\_name}，按层以微米指定）。
\end{itemize}

图~\ref{fig:wire-spreading} 示出 jog 长度与间距的用法。

\begin{figure}[htbp]
\centering
\fbox{\parbox{0.85\textwidth}{\centering\vspace{1.5cm}
\textit{（原书 Figure 90：Wire Spreading Results）}\\[0.3em]
\small 请参阅原版 PDF 插图。
\vspace{1.5cm}}}
\caption{Wire Spreading 结果}
\label{fig:wire-spreading}
\end{figure}

\begin{lstlisting}
fc_shell> spread_wires -min_jog_length 3 \
   -min_jog_spacing_by_layer_name {{M1 0.07} {M2 0.08}}
\end{lstlisting}

疏导后自动进行详细布线迭代修复 DRC。时序保持模式选项：
\opt{-timing\_preserve\_setup\_slack\_threshold}、
\opt{-timing\_preserve\_hold\_slack\_threshold}、
\opt{-timing\_preserve\_nets}。
仅对 slack 不小于阈值的网、\opt{-timing\_preserve\_nets} 所列网，以及同层两 pitch 内的相邻网执行疏导。

\subsection{执行 Wire Widening}
\subsubsection*{Performing Wire Widening}

在详细布线、冗余 via 与 wire spreading 之后，用 \cmd{widen\_wires} 增大平均线宽以减小开路临界面积。默认将所有线加宽到原宽 1.5 倍。可用 \opt{-widen\_widths\_by\_layer\_name} 为每层定义最多五种候选宽度：

\begin{lstlisting}
fc_shell> widen_wires \
   -widen_widths_by_layer_name {{M1 0.07 0.06} {M2 0.08 0.07}}
\end{lstlisting}

加宽会减小邻线间距，可能削弱 spreading 带来的短路临界面积改善。用 \opt{-spreading\_widening\_relative\_weight} 权衡：默认等权；0.0--0.5 偏向 widening（开路）；0.5--1.0 偏向 spreading（短路）。

加宽后进行详细布线迭代修复 DRC；加宽导线不触发 fat wire 间距规则。时序保持选项同 spreading；不对 slack 小于阈值或指定保护网执行加宽。

% ============================================================
\section{插入金属--绝缘体--金属电容}
\subsection*{Inserting Metal-Insulator-Metal Capacitors}

MiM 电容由两层专用导电层夹绝缘体构成，插在两层常规金属之间，见图~\ref{fig:mim-cross-section}。

\begin{figure}[htbp]
\centering
\fbox{\parbox{0.85\textwidth}{\centering\vspace{1.5cm}
\textit{（原书 Figure 91：Cross-Section View of MiM Capacitor Layers）}\\[0.3em]
\small 请参阅原版 PDF 插图（含 M8/MBOT/MTOP/M9 与 viaMimbottom/viaMimtop 等 maskName）。
\vspace{1.5cm}}}
\caption{MiM 电容层剖面示意}
\label{fig:mim-cross-section}
\end{figure}

MiM 电容通常接在电源与地之间，以在电气噪声下维持供电稳定。尺寸常按电源规划的 power strap 几何定制。

用 \cmd{create\_mim\_capacitor\_array} 插入阵列并连接到电源轨。至少指定 MiM 库单元及阵列 $x$/$y$ 增量：

\begin{lstlisting}
fc_shell> create_mim_capacitor_array \
            -lib_cell my_lib/mim_ref_cell \
            -x_increment 20 -y_increment 20
\end{lstlisting}

默认行为：
\begin{itemize}
  \item 对整个 block 插入；可用 \opt{-boundary} 限制到单个矩形/直角多边形区域。插入时不考虑标准单元、macro、placement blockage 或 voltage area；
  \item 方向 R0；可用 \opt{-orientation}（作用于阵列全部单元）；
  \item 命名：\texttt{mimcap!library\_cell\_name!number}；可用 \opt{-prefix}。
\end{itemize}

% ============================================================
\section{插入 Filler Cell}
\subsection*{Inserting Filler Cells}

为保证电源网连通，可在标准单元行空隙中插入 filler；亦常用于加入去耦电容以稳定供电。工具支持单高与多高 filler。

\begin{noteBox}
插入前须用 \cmd{check\_legality} 确认 block 已合法化。
\end{noteBox}

两种主要方法：
\begin{itemize}
  \item \textbf{有序 filler 参考列表}（标准流程）：按列表顺序在每个间隙插入第一个能放入的单元。若工艺要求特定 filler 邻接标准单元，可先插入这些必选 filler，再用标准流程填剩余空隙；
  \item \textbf{阈值电压规则}（VT-based）：按间隙两侧单元的 VT 类型用用户规则选 filler，多用于成熟晶圆厂节点。
\end{itemize}

\subsection{标准 Filler Cell 插入}
\subsubsection*{Standard Filler Cell Insertion}

使用 \cmd{create\_stdcell\_fillers} 插入，用 \cmd{remove\_stdcell\_fillers\_with\_violation} 检查并移除违例 filler。插入时通过布局合法化器遵守先进节点规则与物理约束；若 \appopt{place.legalize.enable\_advanced\_legalizer}=\cmd{true}，使用先进合法化算法做 2D 规则与单元交互检查，可缩短运行时间。

步骤：
\begin{enumerate}
  \item 用 \cmd{check\_legality} 确认合法。注意：\cmd{check\_legality} 假定所有 \texttt{design\_type==filler} 的标准单元都可用于填充，而 \cmd{create\_stdcell\_fillers} 仅插入 \opt{-lib\_cells} 所列单元，二者差异可能导致插入后仍有不可填间隙。
  \item （可选）用 \cmd{create\_left\_right\_filler\_cells} 为特定标准单元两侧插入必选 filler（见「用特定 Filler 邻接标准单元」）。
  \item （可选）\cmd{set\_host\_options -max\_cores n} 启用多核（仅当 technology file 含最小阈值电压、OD 或 TPO 规则时生效）。
  \item 用 \cmd{create\_stdcell\_fillers -lib\_cells} 插入 metal filler（宜从大到小）。可用：
\begin{lstlisting}
fc_shell> get_lib_cells ref_lib/* -filter "design_type==filler"
fc_shell> set FILLER_CELLS \
   [get_object_name [sort_collection -descending \
      [get_lib_cells ref_lib/* -filter "design_type==filler"] area]]
\end{lstlisting}
  \item \cmd{connect\_pg\_net -automatic} 连接 PG。
  \item \cmd{remove\_stdcell\_fillers\_with\_violation} 移除违例 metal filler（可能需多次，因移除会暴露新违例）。
  \item 再插入 nonmetal filler 并连接 PG。
\end{enumerate}

\begin{seeAlsoBox}
Generic ECO Flow for Timing or Functional Changes；移除 Filler Cell。
\end{seeAlsoBox}

\subsection{控制标准 Filler Cell 插入}
\subsubsection*{Controlling Standard Filler Cell Insertion}

默认填充整个 block 水平标准单元行的全部空隙。可控制：
\begin{itemize}
  \item \textbf{按最低漏电选择}：\opt{-leakage\_vt\_order} 按漏电递减给出 VT 层；\opt{-leakage\_vt\_check} 仅检查既有 filler 是否可换成更低漏电单元（最多报告数由 \appopt{chf.create\_stdcell\_fillers.max\_leakage\_vt\_order\_violations} 控制）。须启用 advanced legalizer，且 technology file 中 VT 层 \texttt{maskType=implant}。
  \item \textbf{填充百分比}：\opt{-utilization}；用 \opt{-type\_utilization} 控制各类（如 ULVT/LVT DCAP）比例（总和 $\le$100），可用 \opt{-fill\_remaining} 填剩余，\opt{-prefer\_type\_ordering} 按组内从左到右优先。
  \item \textbf{区域}：\opt{-bboxes}、\opt{-voltage\_area}。
  \item \textbf{电源网违例}：默认不查（省时）；metal filler 建议 \opt{-rules check\_pnet}。
  \item \textbf{防间隙}：\opt{-smallest\_cell\_size} 为 1/2/3 倍 unit site（advanced legalizer 开启时忽略）。
  \item \textbf{Hard placement blockage}：默认遵守；\opt{-ignore\_hard\_blockages} 可忽略。
  \item \textbf{命名}：\texttt{xofiller!...}；\opt{-prefix}/\opt{-separator}。
  \item \textbf{合法化错误}：默认遇错失败；\opt{-continue\_on\_error} 可强制完成。
\end{itemize}

示例：
\begin{lstlisting}
fc_shell> create_stdcell_fillers -lib_cells $FILLER_CELLS \
   -type_utilization { {*/DCAP16*ULVT */DCAP8*ULVT ...} 70
                       {*/DCAP16*LVT */DCAP8*LVT ...} 30 } \
   -prefer_type_ordering -fill_remaining
\end{lstlisting}

\subsection{检查 Filler Cell DRC 违例}
\subsubsection*{Checking for Filler Cell DRC Violations}

默认 \cmd{remove\_stdcell\_fillers\_with\_violation} 检查名称含 \texttt{xofiller} 的实例与顶层信号布线对象之间的 DRC，并移除违例者。可控制：
\begin{itemize}
  \item \opt{-name}：识别 filler 的字符串；
  \item \opt{-boundary}：检查区域；
  \item \opt{-check\_between\_fixed\_objects true}：与所有邻接对象检查；
  \item \opt{-shorts\_only true}：仅因与布线形状短路而移除；
  \item 双重图案化：当 \appopt{route.common.color\_based\_dpt\_flow}=\cmd{true} 且 technology file 有相应规则时检查；
  \item \opt{-check\_only true}：只检查不移除，写 \texttt{block\_fillers\_with\_violation.rpt}。
\end{itemize}

\subsection{修复剩余 Mask 设计规则违例}
\subsubsection*{Fixing Remaining Mask Design Rule Violations}

部分 mask 规则因出现概率低、计算重而未在插入时处理。插入后用 \cmd{replace\_fillers\_by\_rules -replacement\_rule} 修复。表~\ref{tab:replace-fillers-rules} 列出支持的规则与关键字（原书 Table 34）。

\begin{table}[htbp]
\centering
\caption{\cmd{replace\_fillers\_by\_rules} 支持的设计规则（原书 Table 34）}
\label{tab:replace-fillers-rules}
\begin{tabular}{@{}llp{0.42\textwidth}@{}}
\toprule
设计规则 & \texttt{-replacement\_rule} & 说明 \\
\midrule
End orientation & \texttt{end\_orientation} & 翻转最左/最右 filler 方向 \\
Half row adjacency & \texttt{half\_row\_adjacency} & 半高 filler 邻接限制 \\
Illegal abutment & \texttt{illegal\_abutment} & 禁止某些 filler 邻接 \\
Maximum horizontal length & \texttt{od\_horiztonal\_distance} & 连续 filler 行水平长度 \\
Maximum horizontal edge length & \texttt{horizontal\_edges\_distance} & 特定层邻接对象水平边长 \\
Maximum stacking for small fillers & \texttt{small\_filler\_stacking} & 小 filler 堆叠高度 \\
Maximum vertical edge length & \texttt{max\_vertical\_constraint} & 堆叠 filler 垂直边长 \\
Tap cell spacing & \texttt{od\_tap\_distance} & Tap OD 与邻接 OD 距离 \\
Random & \texttt{random} & 随机替换为定制 filler \\
\bottomrule
\end{tabular}
\end{table}

\subsubsection{End Orientation 规则}
\subsubsection*{End Orientation Rule}

对连续水平单元序列，修改最左/最右且属于目标列表的 filler 方向。须指定 \opt{-target\_fillers}，并用 \opt{-left\_end}/\opt{-right\_end}（\texttt{inverse}、\texttt{R0\_or\_MX}、\texttt{MY\_or\_R180}）。\opt{-check\_only} 仅报告将改方向的单元。

\subsubsection{Half Row Adjacency 规则}
\subsubsection*{Half Row Adjacency Rule}

按邻接标准单元类型（Inbound：顶/底边界在半高处，需 ICIP；Regular：全高边界，邻接侧不得有 ICIP）替换半高 filler。替换单元与原 filler 同高、同宽、同 VT；库提供带 ICIP 的 p/n 型。图~\ref{fig:half-height-filler}；表~\ref{tab:half-height-replace}。

\begin{figure}[htbp]
\centering
\fbox{\parbox{0.85\textwidth}{\centering\vspace{1.2cm}
\textit{（原书 Figure 92：Half-Height Filler Cell Replacement With Half Row Adjacency Rule）}\\[0.3em]
\small 请参阅原版 PDF 插图。
\vspace{1.2cm}}}
\caption{半高 Filler 替换与半行邻接规则}
\label{fig:half-height-filler}
\end{figure}

\begin{table}[htbp]
\centering
\caption{半高替换 Filler Cell（原书 Table 35）}
\label{tab:half-height-replace}
\begin{tabular}{@{}ll@{}}
\toprule
邻接 inbound 标准单元 & 指定替换 filler 的选项 \\
\midrule
无 & \opt{-p\_none}、\opt{-n\_none} \\
仅左 & \opt{-p\_left}、\opt{-n\_left} \\
仅右 & \opt{-p\_right}、\opt{-n\_right} \\
左与右 & \opt{-p\_left\_right}、\opt{-n\_left\_right} \\
顶或底（仅 2X filler） & \opt{-p\_left\_center\_right}、\opt{-n\_left\_center\_right} \\
\bottomrule
\end{tabular}
\end{table}

每个列表每种尺寸只能含一个单元。

\subsubsection{Illegal Abutment 规则}
\subsubsection*{Illegal Abutment Rule}

\begin{lstlisting}
fc_shell> replace_fillers_by_rules -replacement_rule illegal_abutment \
   -replace_abutment { {FILL1a CFILL1a} {FILL1b CFILL1b} } \
   -illegal_abutment {FILLx FILLy}
\end{lstlisting}

\opt{-skip\_fixed\_cells} 可忽略 \texttt{physical\_status} 为 fixed/locked 的 filler。图~\ref{fig:illegal-abutment}。

\begin{figure}[htbp]
\centering
\fbox{\parbox{0.85\textwidth}{\centering\vspace{1.2cm}
\textit{（原书 Figure 93：Fixing an Illegal Abutment Violation）}\\[0.3em]
\small 请参阅原版 PDF 插图。
\vspace{1.2cm}}}
\caption{修复非法邻接违例}
\label{fig:illegal-abutment}
\end{figure}

\subsubsection{最大水平边长规则}
\subsubsection*{Maximum Horizontal Edge Length Rule}

\opt{-replacement\_rule horizontal\_edges\_distance}，须指定 \opt{-max\_constraint\_length}、\opt{-layer}，以及 \opt{-refill\_table} 或 \opt{-constraint\_fillers}/\opt{-non\_constraint\_fillers}。

\begin{lstlisting}
fc_shell> replace_fillers_by_rules \
   -replacement_rule horizontal_edges_distance \
   -max_constraint_length 30 -layer M2 \
   -refill_table { {{VT1_FILL1} {VT1_DECAP4 VT1_DECAP2}}
                   {{VT2_FILL1} {VT2_DECAP4 VT2_DECAP2}} }
\end{lstlisting}

\subsubsection{最大水平长度规则}
\subsubsection*{Maximum Horizontal Length Rule}

\opt{-replacement\_rule od\_horiztonal\_distance}（原文关键字拼写如此）。示例限制连续 filler 水平长度 70\,µm：

\begin{lstlisting}
fc_shell> replace_fillers_by_rules \
   -replacement_rule od_horiztonal_distance -max_constraint_length 70 \
   -refill_table { {{VT1_FILL1} {VT1_DECAP4 VT1_DECAP2}}
                   {{VT2_FILL1} {VT2_DECAP4 VT2_DECAP2}} }
\end{lstlisting}

图~\ref{fig:max-vedge}（原书 Figure 94）。

\begin{figure}[htbp]
\centering
\fbox{\parbox{0.85\textwidth}{\centering\vspace{1.2cm}
\textit{（原书 Figure 94：相关最大长度违例修复示意）}\\[0.3em]
\small 请参阅原版 PDF 插图。
\vspace{1.2cm}}}
\caption{修复最大水平长度相关违例}
\label{fig:max-vedge}
\end{figure}

\subsubsection{小 Filler 最大堆叠规则}
\subsubsection*{Maximum Stacking for Small Filler Cells Rule}

\opt{-replacement\_rule small\_filler\_stacking}，指定 \opt{-max\_constraint\_length}、\opt{-small\_fillers}、\opt{-replacement\_fillers}。超出高度时用大 filler 替换相邻 2--3 个小 filler；无法替换则重排。保持违例处 VT，非违例处 VT 可能改变且不检查。

\begin{lstlisting}
fc_shell> replace_fillers_by_rules \
   -replacement_rule small_filler_stacking -max_constraint_length 10 \
   -small_fillers {FILL1 FILL2} -replacement_fillers {FILL3 FILL4}
\end{lstlisting}

图~\ref{fig:max-stack-replace}、\ref{fig:max-stack-reorder}。

\begin{figure}[htbp]
\centering
\fbox{\parbox{0.85\textwidth}{\centering\vspace{1.2cm}
\textit{（原书 Figure 95：Fixing Maximum Stacking Violation by Replacement）}\\[0.3em]
\small 请参阅原版 PDF 插图。
\vspace{1.2cm}}}
\caption{通过替换修复最大堆叠违例}
\label{fig:max-stack-replace}
\end{figure}

\begin{figure}[htbp]
\centering
\fbox{\parbox{0.85\textwidth}{\centering\vspace{1.2cm}
\textit{（原书 Figure 96：Fixing Maximum Stacking Violation by Reordering）}\\[0.3em]
\small 请参阅原版 PDF 插图。
\vspace{1.2cm}}}
\caption{通过重排修复最大堆叠违例}
\label{fig:max-stack-reorder}
\end{figure}

\subsubsection{最大垂直边长规则}
\subsubsection*{Maximum Vertical Edge Length Rule}

\opt{-replacement\_rule max\_vertical\_constraint}。指定 \opt{-max\_constraint\_length}、\opt{-constraint\_fillers}（默认所有标准单元受约束，\opt{-exception\_cells} 排除；\opt{-no\_violation\_along\_exception\_cells}）、以及 \opt{-non\_constraint\_fillers}/\opt{-non\_constraint\_left\_fillers}/\opt{-non\_constraint\_right\_fillers}。

\begin{lstlisting}
fc_shell> replace_fillers_by_rules \
   -replacement_rule max_vertical_constraint -max_constraint_length 180 \
   -constraint_fillers {DECAP8 DECAP4} -non_constraint_fillers {FILL1} \
   -exception_cells {BUF1 TAP1}
\end{lstlisting}

图~\ref{fig:max-vedge2}。

\begin{figure}[htbp]
\centering
\fbox{\parbox{0.85\textwidth}{\centering\vspace{1.2cm}
\textit{（原书 Figure 97：Fixing a Maximum Vertical Edge Length Violation）}\\[0.3em]
\small 请参阅原版 PDF 插图。
\vspace{1.2cm}}}
\caption{修复最大垂直边长违例}
\label{fig:max-vedge2}
\end{figure}

\subsubsection{Tap Cell Spacing 规则}
\subsubsection*{Tap Cell Spacing Rule}

\opt{-replacement\_rule od\_tap\_distance}。指定 \opt{-tap\_distance\_range}（site 数）、\opt{-tap\_cells}，以及 \opt{-left\_violation\_tap}/\opt{-right\_violation\_tap}/\opt{-both\_violation\_tap}。目标与替换 tap 须同尺寸。对无 OD 邻接单元用 \opt{-adjacent\_non\_od\_cells} 跳过。

\begin{lstlisting}
fc_shell> replace_fillers_by_rules -replacement_rule od_tap_distance \
   -tap_cells {TAP} -tap_distance_range {3 3} \
   -left_violation_tap {TAPl} -right_violation_tap {TAPr} \
   -both_violation_tap {TAPb}
\end{lstlisting}

图~\ref{fig:tap-spacing-fix}。

\begin{figure}[htbp]
\centering
\fbox{\parbox{0.85\textwidth}{\centering\vspace{1.2cm}
\textit{（原书 Figure 98：Fixing a Tap Cell Spacing Violation）}\\[0.3em]
\small 请参阅原版 PDF 插图。
\vspace{1.2cm}}}
\caption{修复 Tap Cell 间距违例}
\label{fig:tap-spacing-fix}
\end{figure}

\subsection{随机 Filler Cell 替换}
\subsubsection*{Random Filler Cell Replacement}

\begin{lstlisting}
fc_shell> replace_fillers_by_rules -replacement_rule random \
   -random_replace { {FILL1 {FILL1a FILL1b FILL1c}}
                     {FILL2 {FILL2a FILL2b}} }
\end{lstlisting}

可用 \opt{-skip\_fixed\_cells}。

\subsection{用特定 Filler 邻接标准单元}
\subsubsection*{Abutting Standard Cells With Specific Filler Cells}

用 \cmd{create\_left\_right\_filler\_cells} 在特定标准单元左/右侧插入必选 filler，然后再做标准 filler 插入填剩余空隙。

\opt{-lib\_cells} 规则格式：
\begin{lstlisting}
{standard_cells} {left_filler_cells} {right_filler_cells}
\end{lstlisting}

要求：至少一侧有 filler；filler 的 \texttt{design\_type} 为 \texttt{lib\_cell} 或 \texttt{filler}；与标准单元同 site；标准单元高度为 filler 高度整数倍；宜从大到小。通配符顺序与 \cmd{get\_lib\_cells} 一致。

\begin{lstlisting}
fc_shell> create_left_right_filler_cells \
   -lib_cells {
      { {mylib/SC1} {mylib/LF4X mylib/LF2X} {mylib/RF2X mylib/RF1X} }
      { {mylib/SC2} {mylib/LF2X mylib/LF1X} {} } }
\end{lstlisting}

插入时检查合法位置，但不检查最小 VT、OD、TPO 等先进规则；紧贴标准单元插入，不考虑单元间距规则。

\subsection{控制基于单元的 Filler 插入}
\subsubsection*{Controlling Cell-Based Filler Cell Insertion}

\cmd{create\_left\_right\_filler\_cells} 还可控制：
\begin{itemize}
  \item 区域：\opt{-boundaries}、\opt{-voltage\_areas}；
  \item 方向：默认按行允许方向；\opt{-follow\_stdcell\_orientation} 跟随标准单元（翻转时左右互换）；
  \item 一单位 tile 间隙：\opt{-rules no\_1x}（此时勿在列表中包含宽度为 1 tile 的 filler）；
  \item 命名约定：同标准 filler（\texttt{xofiller!...}）。
\end{itemize}

\subsection{基于阈值电压的 Filler 插入}
\subsubsection*{Threshold-Voltage-Based Filler Cell Insertion}

步骤：
\begin{enumerate}
  \item 确保 VT 类型信息可用；
  \item 用 \cmd{set\_vt\_filler\_rule} 定义规则；
  \item （可选）为多种邻接 VT 组合分别设规则；
  \item 用 \cmd{create\_vtcell\_fillers} 插入。
\end{enumerate}

\begin{lstlisting}
fc_shell> set_vt_filler_rule -vt_type {default default} \
   -filler_cells {myLib/filler2x myLib/filler1x}
fc_shell> set_vt_filler_rule -vt_type {vtA vtA} \
   -filler_cells {myLib/fillerA2x myLib/fillerA1x}
fc_shell> create_vtcell_fillers
\end{lstlisting}

\subsection{控制基于阈值电压的 Filler 插入}
\subsubsection*{Controlling Threshold-Voltage-Based Filler Cell Insertion}

默认填充整个 block 水平标准单元行空隙。可控制 \opt{-region}、\opt{-voltage\_area}，以及命名（默认 \texttt{xofiller!...}；\opt{-prefix}/\opt{-separator}）。

\subsection{移除阈值电压 Filler 信息}
\subsubsection*{Removing the Threshold-Voltage Filler Cell Information}

\begin{lstlisting}
fc_shell> create_vtcell_fillers -clear_vt_information
\end{lstlisting}

\subsubsection{移除 Filler Cell}
\subsubsection*{Removing Filler Cells}

\begin{lstlisting}
fc_shell> remove_cells [get_cells xofiller*]
\end{lstlisting}

仅移除有 DRC 违例的标准流程 filler：用 \cmd{remove\_stdcell\_fillers\_with\_violation}。

% ============================================================
\section{插入 Metal Fill}
\subsection*{Inserting Metal Fill}

布线后，可用金属导线填充 block 空隙以满足金属密度规则。插入前 block 应接近满足时序，且 DRC 很少或没有。

运行 \cmd{signoff\_create\_metal\_fill}，详见第~\ref{chap:icv}~章「使用 IC Validator In-Design 插入 Metal Fill」。

\begin{noteBox}
运行 \cmd{signoff\_create\_metal\_fill} 需要 IC Validator license。
\end{noteBox}

% 章末（原书约 p.544；第 7 章 IC Validator In-Design 之前）
